\documentclass[11pt,a4paper]{article}

% LuaLaTeX ir XeLaTeX naudoja Noto Serif; pdfLaTeX -- numatytąjį šriftą.
\usepackage{iftex}
\ifPDFTeX
  \usepackage[T1]{fontenc}
  \usepackage[utf8]{inputenc}
\else
  \usepackage{fontspec}
  \setmainfont{Noto Serif}
\fi
\usepackage[lithuanian]{babel}
\usepackage{amsmath,amssymb}
\usepackage[a4paper,margin=2.2cm]{geometry}
\usepackage{enumitem}
\usepackage{multicol}
\usepackage{fancyhdr}
\usepackage{microtype}

\setlength{\parindent}{0pt}
\setlength{\parskip}{4pt}
\setlist[enumerate,1]{label=\alph*),leftmargin=2.2em,itemsep=0.45em,topsep=0.4em}
\setlist[enumerate,2]{label=\alph*),leftmargin=2.1em,itemsep=0.2em,topsep=0.3em}
\setlist[itemize]{topsep=0.2em,itemsep=0.1em}
\pagestyle{fancy}
\fancyhf{}
\fancyhead[L]{\small Kvadratiniai daugianariai ir nelygybės}
\fancyhead[R]{\small Užduočių rinkinys}
\fancyfoot[C]{\thepage}
\renewcommand{\headrulewidth}{0.3pt}

\newcounter{uzd}
\newcommand{\uzdavinys}[1][]{%
  \refstepcounter{uzd}%
  \par\medskip\noindent\textbf{\theuzd\ uždavinys\if\relax\detokenize{#1}\relax\else. #1\fi.}\quad
}

\begin{document}

\begin{center}
  {\LARGE\bfseries Kvadratiniai daugianariai ir nelygybės}\par
  \vspace{0.3em}
  {\large Užduočių rinkinys}\par
\end{center}

\noindent\textit{Visuose uždaviniuose laikykite, kad $x\in\mathbb{R}$. Sprendimus pagrįskite.}

\uzdavinys
Mokytojas lentoje užrašė kvadratinį daugianarį
\[
  P(x)=x^2-6x+13, \text{ kur } x \in \mathbb{R}.
\]
ir paklausė Petriuko:

\begin{quote}
--- Kam lygi šio daugianario šaknų suma?

Petriukas, prisiminęs Vijeto teoremą, nedvejodamas atsakė:

--- Šešiems!

Mokytojas atsakė:

--- Sėsk, du!
\end{quote}

Pasipiktinęs Petriukas nusprendė paduoti matematikos mokytoją į teismą. Kokie Petriuko šansai laimėti bylą? Atsakymą pagrįskite matematiškai.

\uzdavinys
Mokytojas pateikė daugianarį
\[
  P(x)=2x^2-10x+8
\]
ir paprašė nustatyti jo šaknų sumą ir sandaugą.

Kaziukas atsakė:
\[
  x_1+x_2=10,\qquad x_1x_2=8.
\]
\begin{quote}
Mokytojas pasakė:

--- Sėsk, du!
\end{quote}

Ar mokytojas sadistas, ar Kaziukas žioplas? Atsakymą pagrįskite.

\uzdavinys
Mokytojas pateikė kvadratinį daugianarį $P(x)$, kurio diskriminantas neigiamas ir kuris turi minimumą. Be to, žinoma, kad $P(2)=P(8)$.

\begin{quote}
Mokytojas paklausė:

--- Kam lygu $P(5)$?

Stasiukas atsakė:

--- Minus keturiems!

Mokytojas pasakė:

--- Sėsk, du!
\end{quote}

Kodėl Stasiukas gavo dvejetą? Atsakymą pagrįskite.

\uzdavinys
Duotas kvadratinis daugianaris
\[
  P(x)=2x^2-4x-6.
\]
\begin{enumerate}
  \item Nustatykite daugianario šaknis.
  \item Užrašykite šį daugianarį dar dviem skirtingais pavidalais.
  \item Nustatykite daugianario ekstremumo tašką, ekstremumą ir jo tipą.
  \item Nustatykite, su kuriomis $x$ reikšmėmis daugianario reikšmės yra:
  \begin{itemize}[label=--,leftmargin=1.5em]
    \item teigiamos;
    \item neigiamos;
    \item neteigiamos;
    \item neneigiamos.
  \end{itemize}
\end{enumerate}

\uzdavinys
Duotas kvadratinis daugianaris
\[
  P(x)=-2(x+1)^2+8.
\]
\begin{enumerate}
  \item Nustatykite daugianario šaknis.
  \item Užrašykite šį daugianarį dar dviem skirtingais pavidalais.
  \item Nustatykite daugianario ekstremumo tašką, ekstremumą ir jo tipą.
  \item Nustatykite, su kuriomis $x$ reikšmėmis daugianario reikšmės yra:
  \begin{itemize}[label=--,leftmargin=1.5em]
    \item teigiamos;
    \item neigiamos;
    \item neteigiamos;
    \item neneigiamos.
  \end{itemize}
\end{enumerate}


\uzdavinys
Duotas kvadratinis daugianaris
\[
  P(x)=2x^2-12x.
\]
\begin{enumerate}
  \item Nustatykite daugianario šaknis.
  \item Užrašykite šį daugianarį dar dviem skirtingais pavidalais.
  \item Nustatykite daugianario ekstremumo tašką, ekstremumą ir jo tipą.
  \item Nustatykite, su kuriomis $x$ reikšmėmis daugianario reikšmės yra:
  \begin{itemize}[label=--,leftmargin=1.5em]
    \item teigiamos;
    \item neigiamos;
    \item neteigiamos;
    \item neneigiamos.
  \end{itemize}
\end{enumerate}

\uzdavinys
Duotas kvadratinis daugianaris
\[
  P(x)=3x^2-12.
\]
\begin{enumerate}
  \item Nustatykite daugianario šaknis.
  \item Užrašykite šį daugianarį dar dviem skirtingais pavidalais.
  \item Nustatykite daugianario ekstremumo tašką, ekstremumą ir jo tipą.
  \item Nustatykite, su kuriomis $x$ reikšmėmis daugianario reikšmės yra:
  \begin{itemize}[label=--,leftmargin=1.5em]
    \item teigiamos;
    \item neigiamos;
    \item neteigiamos;
    \item neneigiamos.
  \end{itemize}
\end{enumerate}

\uzdavinys
Išskaidę kairiąją pusę dauginamaisiais, nustatykite visas realiąsias $x$ reikšmes, su kuriomis teisingos nelygybės:
\begin{multicols}{2}
\begin{enumerate}
  \item $2x^2-2x-12>0$;
  \item $-3x^2+12x-9\geq 0$.
  \item $x^2-5x<0$;
  \item $x^2-9<0$;
  \item $4x^2-25>0$;
  \item $3x ^{2} +2 > 0$ 
  \item $2x^2-4x+2\leq 0$;
  \item $3x^2+12x+12<0$;
  \item $-2x^2+12x-18\geq 0$;
  \item $-3x^2-6x-3>0$.
\end{enumerate}
\end{multicols}

\uzdavinys
Mokytojas pakvietė prie lentos Jonuką rasti visas $x$ reikšmes, su kuriomis nelygybė yra teisinga:
\[
  (x-2)(x+3)>0.
\]
\begin{quote}
Jonukas sako:

--- Easy! Padalijame abi puses iš $(x-2)$, gauname $x+3>0$, taigi $x>-3$.

Mokytojas sako:

--- Sėsk, du!
\end{quote}

Kur problema Jonuko sprendime? Atsakymą pagrįskite ir išspręskite nelygybę teisingai.

\uzdavinys
Mokytojas, pamatęs, kad Jonukas apsiverkė, nusprendė suteikti jam progą pasitaisyti ir ištrinti dvejetą. Tam reikėjo išspręsti lygtį:
\[
  x^2=16.
\]
Jonukas apsidžiaugė, ištraukė kvadratinę šaknį iš abiejų lygties pusių ir parašė:
\[
  x=4.
\]
Mokytojas, susiėmęs už širdies, išėjo pas direktorių ir parašė prašymą išeiti iš darbo...

Kas Jonuko sprendime privertė mokytoją išeiti iš darbo? Atsakymą pagrįskite ir išspręskite lygtį teisingai.

\uzdavinys
Duotas kvadratinis daugianaris
\[
  P(x)=2x^2-8x+11.
\]
Dviem skirtingais būdais įrodykite, kad $P(x)>0$ su visomis realiosiomis $x$ reikšmėmis:
\begin{enumerate}
  \item išskirdami pilnąjį kvadratą;
  \item remdamiesi tik daugianario koeficientais $a,b,c$ ir diskriminantu.
\end{enumerate}

\uzdavinys
Duotas kvadratinis daugianaris
\[
  P(x)=-3x^2+12x-15.
\]
Dviem skirtingais būdais įrodykite, kad $P(x)<0$ su visomis realiosiomis $x$ reikšmėmis:
\begin{enumerate}
  \item išskirdami pilnąjį kvadratą;
  \item remdamiesi tik daugianario koeficientais $a,b,c$ ir diskriminantu.
\end{enumerate}


\uzdavinys
Duotas reiškinys
\[
  T(x)=\sqrt{x^2-5x+4}.
\]
\begin{enumerate}
  \item Apskaičiuokite $T(0)$, $T(2)$, $T(1)$, jeigu šios reikšmės apibrėžtos realiųjų skaičių aibėje.
  \item Nustatykite visas realiąsias $x$ reikšmes, kurias galima įstatyti į reiškinį $T(x)$, kad gautume realųjį skaičių.
\end{enumerate}

\uzdavinys
Duotas reiškinys
\[
  T(x)=\frac{\sqrt{-x^2+6x+7}}{7-x}.
\]
\begin{enumerate}
  \item Apskaičiuokite $T(3)$, $T(7)$, $T(8)$, jeigu šios reikšmės apibrėžtos realiųjų skaičių aibėje.
  \item Nustatykite visas realiąsias $x$ reikšmes, kurias galima įstatyti į reiškinį $T(x)$, kad gautume realųjį skaičių.
\end{enumerate}


\uzdavinys
Duota nelygybė
\[
  3x^2+5x-4>2x^2+4x+2.
\]
\begin{enumerate}
  \item Pertvarkykite nelygybę į pavidalą $ax^2+bx+c>0$.
  \item Išspręskite gautą nelygybę.
\end{enumerate}

\uzdavinys
Duota nelygybė
\[
  (2x+1)^2-9x-6<(1-x)^2.
\]
\begin{enumerate}
  \item Pertvarkykite nelygybę į pavidalą $ax^2+bx+c<0$.
  \item Išspręskite gautą nelygybę.
\end{enumerate}


\uzdavinys
Duotas kvadratinis daugianaris
\[
  P(x)=-2x^2+3x+5
\]
ir tiesinis daugianaris
\[
  Q(x)=x+1.
\]
Nustatykite visas realiąsias $x$ reikšmes, su kuriomis $P(x)>Q(x)$.

\uzdavinys
Duotas kvadratinis daugianaris
\[
  P(x)=2(x+2)(x-1)
\]
ir tiesinis daugianaris
\[
  Q(x)=-2x+2.
\]
Su kuriomis realiosiomis $x$ reikšmėmis daugianario $P(x)$ reikšmės yra \textbf{ne didesnės} už daugianario $Q(x)$ reikšmes?


\uzdavinys
Duoti du kvadratiniai daugianariai:
\[
  P(x)=x^2-4,\qquad Q(x)=x^2-3x.
\]
\begin{enumerate}
  \item Nustatykite intervalus, kuriuose daugianaris $P(x)$ yra teigiamas, ir intervalus, kuriuose jis yra neigiamas.
  \item Nustatykite intervalus, kuriuose daugianaris $Q(x)$ yra teigiamas, ir intervalus, kuriuose jis yra neigiamas.
  \item Raskite bendrus intervalus, kuriuose abu daugianariai yra teigiami.
  \item Raskite bendrus intervalus, kuriuose abu daugianariai yra neigiami.
  \item Raskite visas $x$ reikšmes, su kuriomis abu daugianariai yra neneigiami.
  \item Raskite visas $x$ reikšmes, su kuriomis abu daugianariai yra neteigiami.
\end{enumerate}


\uzdavinys
Duoti du kvadratiniai daugianariai:
\[
  P(x)=-x^2+x+6,\qquad Q(x)=x^2-4.
\]
\begin{enumerate}
  \item Nustatykite intervalus, kuriuose daugianaris $P(x)$ yra teigiamas, ir intervalus, kuriuose jis yra neigiamas.
  \item Nustatykite intervalus, kuriuose daugianaris $Q(x)$ yra teigiamas, ir intervalus, kuriuose jis yra neigiamas.
  \item Raskite bendrus intervalus, kuriuose abu daugianariai yra teigiami.
  \item Raskite bendrus intervalus, kuriuose abu daugianariai yra neigiami.
  \item Raskite visas $x$ reikšmes, su kuriomis abu daugianariai yra neneigiami.
  \item Raskite visas $x$ reikšmes, su kuriomis abu daugianariai yra neteigiami.
  \item Raskite bendrus intervalus, kuriuose $P(x)>0$ ir $Q(x)<0$.
  \item Raskite bendrus intervalus, kuriuose $P(x)<0$ ir $Q(x)>0$.
\end{enumerate}
c)--h) dalyse raskite atitinkamų sprendinių aibių sankirtą.

\uzdavinys
Kvadratinis daugianaris $P(x)=ax^2+bx+c$ nagrinėjamas uždarame intervale $x\in[1;10]$. Žinoma, kad $P(2)=P(6)$ ir $b>0$. Nustatykite, su kuria $x$ reikšme šiame intervale $P(x)$ įgyja mažiausią reikšmę ir su kuria -- didžiausią.

\uzdavinys
Kvadratinis daugianaris $P(x)=ax^2+bx+c$ nagrinėjamas uždarame intervale $x\in[0;4]$. Žinoma, kad $P(3)=P(7)$ ir $b<0$. Nustatykite, su kuria $x$ reikšme šiame intervale $P(x)$ įgyja mažiausią reikšmę ir su kuria -- didžiausią.

\uzdavinys
Kvadratinis daugianaris $P(x)=ax^2+bx+c$ nagrinėjamas uždarame intervale $x\in[1;9]$. Žinoma, kad $P(3)=P(7)$, diskriminantas $D>0$, o ekstremumas lygus $\pi$. Nustatykite visas $x$ reikšmes, su kuriomis šiame intervale $P(x)$ įgyja mažiausią reikšmę, ir visas $x$ reikšmes, su kuriomis įgyja didžiausią.

\uzdavinys
Sukurkite kvadratinį daugianarį $P(x)$, kuris tenkina abi sąlygas:
\begin{enumerate}
  \item $P(x)<0$ tada ir tik tada, kai $1<x<5$;
  \item $P(0)=10$.
\end{enumerate}
Užrašykite gautą daugianarį pavidalu $P(x)=ax^2+bx+c$.

\uzdavinys
Sukurkite kvadratinį daugianarį $P(x)$, kuris tenkina abi sąlygas:
\begin{enumerate}
  \item $P(x)>0$ tada ir tik tada, kai $x<-2$ arba $x>4$;
  \item daugianario ekstremumas lygus $-18$.
\end{enumerate}
Užrašykite gautą daugianarį pavidalu $P(x)=ax^2+bx+c$.

\uzdavinys
Sukurkite kvadratinį daugianarį
\[
  P(x)=ax^2+bx+c,
\]
kuris tenkina abi sąlygas:
\begin{enumerate}
  \item $P(x)>0$ su visomis realiosiomis $x$ reikšmėmis, išskyrus $x=4$;
  \item $P(1)=18$.
\end{enumerate}

\clearpage
\fancyhead[R]{\small Atsakymai}
\section*{Atsakymai}
\noindent\textit{Pateikiami galutiniai atsakymai ir trumpi pagrindimai ten, kur reikia įrodyti teiginį.}

\begin{enumerate}[label=\textbf{\arabic*.},leftmargin=2.6em,itemsep=0.85em,topsep=0.55em]

\item $D=(-6)^2-4\cdot 1\cdot13=-16<0$, todėl realiųjų šaknų nėra ir jų sumos apibrėžti negalima. Sąlygoje aiškiai nurodyta, kad $x\in\mathbb{R}$, taigi kalbama tik apie realiąsias šaknis. Petriuko atsakymas neteisingas, todėl laimėti bylą jis neturi jokių šansų.

\item Kaziukas suklydo: taikydamas Vijeto teoremą, nepadalijo iš vyriausiojo koeficiento $a=2$. Kadangi $D=(-10)^2-4\cdot2\cdot8=36>0$, yra dvi realiosios šaknys. Jų suma $x_1+x_2=-\frac{b}{a}=\frac{10}{2}=5$, o sandauga $x_1x_2=\frac{c}{a}=\frac{8}{2}=4$. Taigi mokytojas pagrįstai nurodė klaidą; taigi mokytojas nėra sadistas.

\item Iš $P(2)=P(8)$ gauname, kad ekstremumo taškas yra $x=\frac{2+8}{2}=5$. Kadangi daugianaris turi minimumą, tai $a>0$. Todėl
  \[
    P(5)=-\frac{D}{4a}>0,
  \]
  nes $D<0$. Vadinasi, $P(5)$ negali būti lygu $-4$, todėl Stasiuko atsakymas neteisingas. Tikslios $P(5)$ reikšmės iš pateiktų duomenų nustatyti negalima.

\item
\begin{enumerate}[label=\alph*),leftmargin=2em,itemsep=0.2em,topsep=0.2em]
  \item $x_1=-1$, $x_2=3$.
  \item $P(x)=2(x+1)(x-3)=2(x-1)^2-8$.
  \item Ekstremumo taškas $x=1$; minimumas $P_{\min}=-8$.
  \item $P(x)>0$, kai $x\in(-\infty;-1)\cup(3;+\infty)$;\\
        $P(x)<0$, kai $x\in(-1;3)$;\\
        $P(x)\le0$, kai $x\in[-1;3]$;\\
        $P(x)\ge0$, kai $x\in(-\infty;-1]\cup[3;+\infty)$.
\end{enumerate}

\item
\begin{enumerate}[label=\alph*),leftmargin=2em,itemsep=0.2em,topsep=0.2em]
  \item $x_1=-3$, $x_2=1$.
  \item $P(x)=-2x^2-4x+6=-2(x+3)(x-1)$.
  \item Ekstremumo taškas $x=-1$; maksimumas $P_{\max}=8$.
  \item $P(x)>0$, kai $x\in(-3;1)$;\\
        $P(x)<0$, kai $x\in(-\infty;-3)\cup(1;+\infty)$;\\
        $P(x)\le0$, kai $x\in(-\infty;-3]\cup[1;+\infty)$;\\
        $P(x)\ge0$, kai $x\in[-3;1]$.
\end{enumerate}

\item
\begin{enumerate}[label=\alph*),leftmargin=2em,itemsep=0.2em,topsep=0.2em]
  \item $x_1=0$, $x_2=6$.
  \item $P(x)=2x(x-6)=2(x-3)^2-18$.
  \item Ekstremumo taškas $x=3$; minimumas $-18$.
  \item $P(x)>0$, kai $x\in(-\infty;0)\cup(6;+\infty)$;\\
        $P(x)<0$, kai $x\in(0;6)$;\\
        $P(x)\le0$, kai $x\in[0;6]$;\\
        $P(x)\ge0$, kai $x\in(-\infty;0]\cup[6;+\infty)$.
\end{enumerate}

\item
\begin{enumerate}[label=\alph*),leftmargin=2em,itemsep=0.2em,topsep=0.2em]
  \item $x_1=-2$, $x_2=2$.
  \item $P(x)=3(x-2)(x+2)=3(x-0)^2-12$.
  \item Ekstremumo taškas $x=0$; minimumas $-12$.
  \item $P(x)>0$, kai $x\in(-\infty;-2)\cup(2;+\infty)$;\\
        $P(x)<0$, kai $x\in(-2;2)$;\\
        $P(x)\le0$, kai $x\in[-2;2]$;\\
        $P(x)\ge0$, kai $x\in(-\infty;-2]\cup[2;+\infty)$.
\end{enumerate}

\item
\begin{enumerate}[label=\alph*),leftmargin=2em,itemsep=0.2em,topsep=0.2em]
  \item $x\in(-\infty;-2)\cup(3;+\infty)$;
  \item $x\in[1;3]$;
  \item $x\in(0;5)$;
  \item $x\in(-3;3)$;
  \item $x\in\left(-\infty;-\frac52\right)\cup\left(\frac52;+\infty\right)$;
  \item $x\in\mathbb{R}$;
  \item $x=1$;
  \item Sprendinių nėra ($\varnothing$);
  \item $x=3$;
  \item Sprendinių nėra ($\varnothing$).
\end{enumerate}

\item Jonukas dalijo iš reiškinio $x-2$, neatsižvelgęs į jo ženklą. Kai $x>2$, dalijame iš teigiamo skaičiaus, todėl nelygybės ženklas nesikeičia: gauname $x+3>0$, o kartu su $x>2$ tai duoda $x>2$. Kai $x<2$, dalijame iš neigiamo skaičiaus, todėl ženklą būtina apversti: gauname $x+3<0$, taigi $x<-3$. Kai $x=2$, dalyti negalima, o pradinė nelygybė neteisinga, nes sandauga lygi nuliui. Atsakymas:
  \[
    x\in(-\infty;-3)\cup(2;+\infty).
  \]

\item Yra du skaičiai, kuriuos pakėlę kvadratu gauname $16$: $4$ ir $-4$. Sprendimas turėjo būti $x=\pm\sqrt{16}=\pm4$. Jonukas prarado sprendinį $x=-4$.

\item a) $P(x)=2(x-2)^2+3\ge 3>0$.\\
      b) $a=2>0$, todėl daugianaris turi minimumą. Apskaičiuojame $D=(-8)^2-4\cdot2\cdot11=-24$. Minimumas lygus
      \[
        -\frac{D}{4a}=-\frac{-24}{4\cdot2}=3.
      \]
      Kadangi $3>0$, tai $P(x)\geq 3>0$ visiems $x\in\mathbb{R}$.

\item a) $P(x)=-3(x-2)^2-3\le -3<0$.\\
      b) $a=-3<0$, todėl daugianaris turi maksimumą. Apskaičiuojame $D=12^2-4\cdot(-3)\cdot(-15)=-36$. Maksimumas lygus
      \[
        -\frac{D}{4a}=-\frac{-36}{4\cdot(-3)}=-3.
      \]
      Kadangi $-3<0$, tai $P(x)\leq -3<0$ visiems $x\in\mathbb{R}$.

\item a) $T(0)=2$; $T(2)$ neapibrėžta realiųjų skaičių aibėje; $T(1)=0$.\\
      b) $x\in(-\infty;1]\cup[4;+\infty)$.

\item a) $T(3)=1$; $T(7)$ ir $T(8)$ neapibrėžtos realiųjų skaičių aibėje.\\
      b) $x\in[-1;7)$.

\item a) $x^2+x-6>0$.\quad b) $x\in(-\infty;-3)\cup(2;+\infty)$.

\item a) $3x^2-3x-6<0$.\quad b) $x\in(-1;2)$.

\item $x\in(-1;2)$.

\item $x\in[-3;1]$.

\item
\begin{enumerate}[label=\alph*),leftmargin=2em,itemsep=0.2em,topsep=0.2em]
  \item $P(x)>0$, kai $x\in(-\infty;-2)\cup(2;+\infty)$; $P(x)<0$, kai $x\in(-2;2)$.
  \item $Q(x)>0$, kai $x\in(-\infty;0)\cup(3;+\infty)$; $Q(x)<0$, kai $x\in(0;3)$.
  \item $x\in(-\infty;-2)\cup(3;+\infty)$;
  \item $x\in(0;2)$;
  \item $x\in(-\infty;-2]\cup[3;+\infty)$;
  \item $x\in[0;2]$.
\end{enumerate}

\item
\begin{enumerate}[label=\alph*),leftmargin=2em,itemsep=0.2em,topsep=0.2em]
  \item $P(x)>0$, kai $x\in(-2;3)$; $P(x)<0$, kai $x\in(-\infty;-2)\cup(3;+\infty)$.
  \item $Q(x)>0$, kai $x\in(-\infty;-2)\cup(2;+\infty)$; $Q(x)<0$, kai $x\in(-2;2)$.
  \item $x\in(2;3)$;
  \item Sprendinių nėra ($\varnothing$);
  \item $x\in\{-2\}\cup[2;3]$;
  \item $x=-2$;
  \item $x\in(-2;2)$;
  \item $x\in(-\infty;-2)\cup(3;+\infty)$.
\end{enumerate}

\item Iš $P(2)=P(6)$ gauname, kad ekstremumo taškas yra $x=\frac{2+6}{2}=4$. Kadangi $-\frac{b}{2a}=4$, tai $b=-8a$. Iš $b>0$ gauname $a<0$, todėl didžiausia reikšmė įgyjama, kai $x=4$. Kuo toliau nuo $4$, tuo daugianario reikšmė mažesnė. Intervalo galų atstumai nuo $4$ yra $4-1=3$ ir $10-4=6$. Tolimiausias intervalo skaičius yra $10$, todėl mažiausia reikšmė įgyjama, kai $x=10$.

\item Iš $P(3)=P(7)$ gauname, kad ekstremumo taškas yra $x=\frac{3+7}{2}=5$. Jis nepriklauso intervalui $[0;4]$. Kadangi $-\frac{b}{2a}=5$, tai $b=-10a$. Iš $b<0$ gauname $a>0$, todėl kuo toliau nuo $5$, tuo daugianario reikšmė didesnė. Šiame intervale arčiausiai $5$ yra skaičius $4$ (atstumas $1$), o toliausiai -- skaičius $0$ (atstumas $5$). Todėl mažiausia reikšmė intervale įgyjama, kai $x=4$, o didžiausia -- kai $x=0$.

\item Iš $P(3)=P(7)$ gauname, kad ekstremumo taškas yra $x=5$. Ekstremumas lygus $-\frac{D}{4a}=\pi>0$. Kadangi $D>0$, tai $a<0$, todėl ekstremumas yra maksimumas. Didžiausia reikšmė $\pi$ įgyjama, kai $x=5$. Kuo toliau nuo $5$, tuo daugianario reikšmė mažesnė. Abu intervalo galai nuo $5$ nutolę po $4$ vienetus ir yra tolimiausi intervalo skaičiai. Todėl mažiausia reikšmė įgyjama, kai $x=1$ ir kai $x=9$.

\item $P(x)=2(x-1)(x-5)=\boxed{2x^2-12x+10}$.

\item $P(x)=2(x+2)(x-4)=\boxed{2x^2-4x-16}$.

\item $P(x)=2(x-4)^2=\boxed{2x^2-16x+32}$.
\end{enumerate}

\end{document}
